% TOP_PROBLEM: {"id": 30000030, "title": "Cannon's Conjecture", "category_id": 6, "category": {"id": 6, "name": "geometry", "display_name": "Geometry", "description": "Euclidean and non-Euclidean geometry, geometric structures.", "slug": "geometry", "order_index": 6, "created_at": "2026-07-31T15:26:25.671Z"}, "status": "open"}
% FIRST_POSTED: 2026-09-25
% ENTRY_KIND: statement_only
% !TeX program = lualatex
% Definition and sources only; this file is not an LLM solution attempt.
\documentclass[11pt]{article}
\usepackage[margin=25mm]{geometry}
\usepackage{fontspec}
\setmainfont{Latin Modern Roman}
\usepackage{amsmath,amssymb,mathtools}
\usepackage{unicode-math}
\setmathfont{Latin Modern Math}
\usepackage{xurl,hyperref}
\hypersetup{colorlinks=true,urlcolor=blue}
\setcounter{secnumdepth}{0}
\setlength{\parindent}{0pt}
\setlength{\parskip}{5pt}
\setlength{\emergencystretch}{3em}
\begin{document}
\section*{\texorpdfstring{Cannon's Conjecture}{Cannon's Conjecture}}\label{30000030}
\subsection{Definitions and mathematical statement}
A finitely generated group is Gromov-hyperbolic if a Cayley graph has uniformly thin geodesic triangles. Its boundary consists of asymptotic classes of geodesic rays, with the standard visual topology. The conjecture is
\[
 \partial G\cong S^2\quad\Longrightarrow\quad
 G\curvearrowright\mathbb H^3\text{ properly, cocompactly, and isometrically}.
\]
Properness allows finite stabilizers in the presence of torsion. Equivalently the group is virtually a cocompact Kleinian group in the source's convention. The hypothesis is about the boundary's topology, not merely its dimension.
\par\smallskip\textit{Formulation and concept references:} \href{https://www.ams.org/bookstore/pspdf/surv-225-prev.pdf}{[S1]}.

\subsection{Short English statement}
Must a hyperbolic group with a two-sphere boundary come from a cocompact hyperbolic three-dimensional geometry?
\subsection{Sources}
\begin{itemize}
\item[S1]Expanding Thurston Maps.
\url{https://www.ams.org/bookstore/pspdf/surv-225-prev.pdf}.
\textit{Formulation source.}
\end{itemize}
\end{document}
